% ============================================================
%  Paper 17: Inflationary Observables from Dimensional Flow in Scale-Dependent Geometric Field Theory
%  Target: JCAP Compilation (SDGFT Series)
% ============================================================
\documentclass[a4paper,11pt]{article}
\usepackage{jcappub}

\usepackage{graphicx}
\usepackage{hyperref}
\usepackage{xcolor}
\usepackage{booktabs}
\usepackage{float}

% ── SDGFT shared macros ──────────────────────────────────────
\usepackage{sdgft-macros}

% ── Document ─────────────────────────────────────────────────
\begin{document}

\title{Inflationary Observables from Dimensional Flow in Scale-Dependent Geometric Field Theory}

\author{David A. Besemer}
\affiliation{Independent Researcher}
\emailAdd{david.besemer@sdgft.org}

\abstract{
We derive the inflationary observables of SDGFT, where inflation is driven by the dimensional flow from the UV fixed point to the IR attractor. We calculate the primordial spectral index ns ~ 0.967 and the tensor-to-scalar ratio r ~ 0.013. We present the formal erratum correcting the legacy blue spectrum formula.
}

\arxivnumber{SDGFT-2026-17}

\maketitle

\section{Introduction}
Inflation is typically driven by an inflaton field. In SDGFT, it arises from the transition of the spectral dimension of spacetime.

\section{Theoretical Formulation}
Here we formulate the core mathematical relations governing the system:
\begin{equation}
  \ns = 1 - \frac{2(2n-1)}{\Ne(2n-1) + n} \approx 0.9671
\end{equation}
\begin{equation}
  \rts = \frac{48(2n-1)^2}{[\Ne(2n-1) + n]^2} \approx 0.0130
\end{equation}
\section{CMB Observables}
The scalar spectral index and tensor-to-scalar ratio are calculated using the Jordan-frame slow-roll equations for f(R) gravity.

\section{Conclusion}
The predicted tensor-to-scalar ratio r ~ 0.013 will be tested by next-generation polarization experiments, providing a test of the theory.



\begin{thebibliography}{99}
\bibitem{Goldstein_Paper01} D. A. Besemer, ``Axiomatic Foundations of SDGFT,'' SDGFT-2026-01 (in preparation).
\bibitem{Goldstein_Paper04} D. A. Besemer, ``Effective Spectral Dimension and the Banach Fixed-Point Theorem in SDGFT,'' SDGFT-2026-04.
\bibitem{Goldstein_Code} D. A. Besemer, \texttt{sdgft} Python package, \url{https://github.com/david-besemer/sdgft} (2026).
\end{thebibliography}

\end{document}
